\documentclass[11pt]{article}

\usepackage[margin=1in]{geometry}
\usepackage{amsmath,amssymb}
\usepackage{booktabs}
\usepackage{enumitem}
\usepackage[hidelinks]{hyperref}
\usepackage{xurl}
\usepackage{xcolor}

\setlength{\parindent}{0pt}
\setlength{\parskip}{0.65em}
\setlist[itemize]{leftmargin=1.5em,itemsep=0.25em,topsep=0.25em}

\title{\vspace{-2em}Uniform(1,2) Canonical Simulation Update}
\author{}
\date{\vspace{-2em}}

\begin{document}
\maketitle

\section*{Purpose}

This simulation restarts the fixed-DGP IFEval-like comparison using weaker
nonzero loadings and an apples-to-apples canonical parameterization. The earlier
Uniform(2,3) loading design was useful as a high-signal case, but it can make
binary responses nearly deterministic for low and high latent classes. The new
run uses milder loadings to diagnose whether factor-score collapse and
parameter discrepancies are driven by saturation of the binary observations.

\section*{Simulation Design}

For each design cell, the structural parameters are fixed and a new binary data
set is simulated in each replication. The loading matrix uses the IFEval-like
unbalanced block allocation, with at least 30 primary items in the smallest
block. Primary and cross-loading magnitudes are both sampled from
\(\mathrm{Uniform}(1,2)\), cross-loadings occur with probability 0.05, and
cross-loading signs are randomly assigned.

\begin{center}
\begin{tabular}{ll}
\toprule
Dimension & Values \\
\midrule
Sample size & \(n \in \{100,200,400\}\) \\
Items, Product MAP & \(p \in \{500,1000,1500,2000\}\) \\
Items, Gibbs baselines & \(p \in \{500,1000\}\) \\
Latent factors & \(H \in \{5,10\}\) \\
Mixture components & \(G_h = 2\) for all \(h\), or \(G_h = 3\) for all \(h\) \\
Mixture separation & \(\mathrm{sep}=2\) \\
Replications & 25 per setting \\
\bottomrule
\end{tabular}
\end{center}

Product MAP and Viroli Laplace Gibbs use the same loading shrinkage level in
matched cells:
\[
\lambda_{\Lambda}(n) =
\begin{cases}
3, & n \in \{100,200\},\\
6, & n = 400.
\end{cases}
\]
The diffuse Gaussian Viroli baseline has no Laplace loading penalty.

\section*{Canonicalization}

All reported parameter RMSEs use the same canonical factor scale. For each
factor coordinate, the fitted marginal mixture is transformed to have mean zero
and variance one. If
\[
\tilde f_{ih} = \frac{f_{ih} - m_h}{s_h},
\]
then loadings and intercepts are transformed so that the probit linear
predictor is unchanged:
\[
\alpha_j + \sum_{h=1}^H f_{ih}\lambda_{jh}
=
\tilde \alpha_j + \sum_{h=1}^H \tilde f_{ih}\tilde \lambda_{jh}.
\]
Product MAP is canonicalized after fitting before evaluation. Viroli Gibbs
normalizes each retained posterior draw using the same convention before
posterior averaging. This makes factor, loading, intercept, mixture-mean, and
mixture-variance RMSEs comparable across methods.

\section*{Computation Order}

The run is staged deliberately.

\begin{enumerate}
\item Product MAP is run first for the full \(p\)-grid. Replications are
      launched serially, while each fit uses 18 internal workers for
      conditionally independent updates.
\item After Product MAP completes, the Gibbs baselines are run for
      \(p \in \{500,1000\}\). Independent Gibbs jobs are launched in parallel,
      and each Gibbs fit uses 4 internal workers for conditionally independent
      blocks. The Markov chain iterations themselves remain serial.
\end{enumerate}

This separates the fast deterministic estimator from the more expensive Gibbs
baselines and avoids oversubscribing the machine while Product MAP is using the
full worker pool.

\section*{Recorded Outputs}

The combined results file records timing, convergence flags, factor RMSE,
probability-score RMSE, intercept RMSE, loading RMSE, mixture mean RMSE,
mixture variance RMSE, mixture weight RMSE, and Gibbs effective sample-size
summaries when available.

The run can be reproduced from the repository root with:
\begin{verbatim}
zsh scripts/sample_size/run_uniform12_canonical_product_first_simulation.sh
\end{verbatim}

The active output directory is:
\begingroup
\footnotesize
\path{results/full/fixed_ifeval_lambda_min30_u1_2_cp0_05_sep2_npenalty3_6_h5_h10_canonical_productfirst_gibbsparallel/}
\par
\endgroup

The final combined CSV will be:
\begingroup
\footnotesize
\path{results/full/fixed_ifeval_lambda_min30_u1_2_cp0_05_sep2_npenalty3_6_h5_h10_canonical_productfirst_gibbsparallel/comparison_results.csv}
\par
\endgroup

\end{document}
